% FILE: 06_contribution_to_thesis.tex
% Project: audits
% Role: completeness-audit-and-compliance-verification-corpus

\section{Contribution to the Dissertation}
\label{sec:audits-contribution}

\subsection{Contribution to Prediction vs.\ Coordination}
\label{sec:audits-contribution-prediction}

The central thesis distinguishes \emph{prediction}~--- inference from historical
patterns~--- from \emph{coordination}~--- real-time synchronization of process
state across agents and artifacts under cryptographic receipt discipline.

The audits corpus contributes to this distinction by demonstrating that quality
enforcement requires coordination, not prediction. A machine-learning model
trained on historical test outcomes could predict, with high probability, that
the \texttt{wasm4pm-compat} crate is conformant: after all, 4 of 5 audits pass,
84 tests pass, and the forensic audit returns \textsc{clean}. But the prediction
would be wrong. The \textsc{dto\_001} violation at \texttt{mod.rs:735} is
undetectable by behavioral testing; it is detectable only by a deterministic
pattern audit that queries the type-law surface of the source artifact.

Coordination, in this context, means that the audit must run, the receipt must be
committed, and the graduation checkpoint must be withheld until the receipt is
\textsc{pass}. No amount of prediction can substitute for this causal chain of
evidence. The audits corpus is the concrete proof that the coordination model
catches defects the prediction model cannot.

The EWMA concept drift detection mechanism contributes a complementary result:
statistical process control (a predictive technique) is insufficient for
conformance monitoring when the underlying process law has changed. Within 4
traces of injected drift, the EWMA statistic crosses the \(LCL = 0.92\) boundary,
triggering a coordination-level alert that stops the process and demands human
review. The alert is a coordination event, not a prediction.

\subsection{Contribution to Process Evidence}
\label{sec:audits-contribution-evidence}

The dissertation argues that process evidence must satisfy three properties: it
must be derivable from event logs (not from code assertions), it must be auditable
against a declared process model, and it must carry a cryptographic receipt that
can be replayed.

The audits corpus demonstrates all three properties at the system level. First,
the five deterministic shell audits are event-log derivations in miniature: they
extract the type-law surface of the codebase as a set of observable events
(method names, feature dependencies, public item counts) and compare these events
against the declared process model (the graduation boundary specification).
Second, the comparison is explicit and reproducible: the declared graduation
boundary is \texttt{compilation \&\& audit\_pass}, and the actual surface is
verified to be a conformant superset of the declared 6-item manifest (87 actual
public items versus 6 declared, with the superset lawfully encompassing the
declared minimum). Third, the receipt chain~--- five committed text files plus
four agent SHA-256 signatures plus the BLAKE3+Ed25519 \texttt{ALIVE\_001}
certification~--- satisfies the cryptographic replay requirement.

The compile-fail fixture corpus contributes an additional evidence layer: 199
type-law receipts that prove, by Rust's type checker, that 97.7\% of the
specified violations are correctly rejected at compile time. This transforms
type-law claims from documentation assertions into machine-verified process
evidence.

\subsection{Contribution to Manufacturing Architecture}
\label{sec:audits-contribution-manufacturing}

The audits corpus contributes to the manufacturing architecture by establishing
the \emph{audit layer} as a first-class manufacturing stage. In the CodeManufactory
doctrine, manufacturing stages are not merely sequential steps: they are lawful
objects with defined receipt shapes, proof gates, and admission refusal
signatures. The audits corpus instantiates this doctrine for the quality
enforcement dimension.

Concretely, the five-audit program defines a manufacturing stage with:
\begin{itemize}
  \item \textbf{Input:} The \texttt{wasm4pm-compat} source tree.
  \item \textbf{Operators:} Five deterministic shell scripts.
  \item \textbf{Proof gate:} All five scripts return \textsc{pass}.
  \item \textbf{Receipt:} Five committed \texttt{.txt} files with
    re-executable audit scripts.
  \item \textbf{Admission refusal:} Any \textsc{fail} result blocks the
    downstream graduation checkpoint.
  \item \textbf{Artifact:} The
    \texttt{PROCESS\_INTELLIGENCE\_COMPAT\_CONFORMANCE\_001} checkpoint file.
\end{itemize}

This pattern~--- operator, proof gate, receipt, refusal, artifact~--- is the
canonical manufacturing stage structure. The audits corpus provides the only
instantiation of this structure that is itself subject to Van der Aalst
conformance auditing, making it the self-validating layer of the manufacturing
architecture.

\subsection{Contribution to Capital-Flow Infrastructure}
\label{sec:audits-contribution-capital}

The board claim support audit (\texttt{audit-board-claim-support.md}) directly
connects the audits corpus to the capital-flow and settlement infrastructure of
the dissertation's final chapter. Board claims in M\&A contexts~--- EBITDA
optimization, working capital improvement, GRC defensibility, integration
velocity~--- must be traceable to process-intelligence evidence if they are to be
board-admissible.

The audit confirms traceability for all nine claim types across four strategic
domains. Each claim type has a verified research support citation (academic paper
and formula), a receipt shape specification (schema, required fields, signing
protocol), a public standard reference (IEEE/OCEL/BPMN), and a lifecycle state
mapping (Autonomic MAPE-K stage). This four-layer traceability chain is the
minimum structure required for a claim to be admitted to a board process.

The audit thereby establishes the architectural bridge between the technical
manufacturing infrastructure (type-law receipts, cryptographic ledgers, Petri net
soundness verification) and the capital-allocation decisions that motivate the
enterprise's investment in process intelligence. Claims that cannot be traced
through this bridge are refused as board-inadmissible~--- applying the same
admission refusal discipline at the business layer that the \texttt{Evidence<T>}
type applies at the code layer.
